\documentclass[12pt]{article}
\usepackage[T1]{fontenc}
\usepackage[utf8]{inputenc}
\usepackage{lmodern}
\usepackage{textcomp}
\usepackage{enumitem}
\usepackage{geometry}
\usepackage{graphicx}
\usepackage{amsmath, amssymb}
\geometry{margin=1in}
\begin{document}
\begin{center}
Political Science - Undergraduate Level Assessment 10 \
Total Marks: 40 \
Answer all questions.
\end{center}

\section*{Multiple Choice (10 marks)}
\begin{enumerate}[label=\arabic*.]
\item Of the following interest groups, which has created the largest number of political action committees (PACs) since the 1970 s?
\begin{enumerate}[label=(\alph*)]
    \item Environmental activists
    \item Labor unions
    \item For-profit business
    \item Religious institutions
\end{enumerate}
\item How did relations with Russia develop under George H.W. Bush?
\begin{enumerate}[label=(\alph*)]
    \item Cautious support for Gorbachev and Yeltsin
    \item Unreserved support for Gorbachev and Yeltsin
    \item Denouncement of Gorbachev, but support for Yeltsin
    \item Denouncement of Gorbachev and Yeltsin
\end{enumerate}
\item In what way did the George W Bush administration change the direction of US foreign policy?
\begin{enumerate}[label=(\alph*)]
    \item It criticized international organizations, rather than trying to strengthen them
    \item It expanded NATO to include former Soviet states
    \item It focused on a more personal style of leadership
    \item It increased international support for the United States
\end{enumerate}
\item Bureaucratic politics suggests we should be worried about which of the following with regard to nuclear weapons?
\begin{enumerate}[label=(\alph*)]
    \item Having the capability to deter the most powerful rival
    \item Having the capability to deter smaller states
    \item How nuclear attacks are identified and responded to; who controls the weapons
    \item Bureaucratic politics provides no information about nuclear proliferation and use
\end{enumerate}
\item According to Realists, what accounts for the onset of the Cold War?
\begin{enumerate}[label=(\alph*)]
    \item Ideological differences
    \item A power vacuum
    \item The expansionist nature of the Soviet Union
    \item Both b and c
\end{enumerate}
\end{enumerate}

\section*{True / False (10 marks)}
\textit{Answer True or False}
\begin{enumerate}[label=\arabic*.]
\setcounter{enumi}{5}
\item True or False: Krauthammer's concept suggested that the post-Cold War system faced no significant threats from other nations.
\item True or False: The Human Development Report was first published in 1987, setting the foundation for development studies.
\item True or False: In 'Federalist No. 10,' James Madison argues that a federal system makes it difficult for one faction to gain governing power.
\item True or False: Intergovernmental organizations often provide robust enforcement to ensure adherence to international agreements.
\item True or False: The National Security Council is the government body responsible for coordinating American foreign and military policy.
\end{enumerate}

\section*{Long Form Response (20 marks)}
\textit{Answer each question with a clear and complete explanation.}
\begin{enumerate}[label=\arabic*.]
\setcounter{enumi}{10}
\item Discuss the various challenges facing security studies in the contemporary world and analyze how these challenges may shape the future of the field.
\item Discuss the implications of using military force to expel social groups from a state, considering its impact on national identity, human rights, and international relations.
\end{enumerate}

\vfill
\noindent\textit{End of Paper}
\end{document}